\section{The Complex Exponential}
\label{sec:esponenziale_complesso}%

The real exponential and trigonometric functions can be thought of as
intermediate tools to define a natural isomorphism from $\CC$ as an additive
group to $\CC$ as a multiplicative group.

Indeed, we can define the function $\exp \colon \CC \to \CC$ via the
\emph{Euler's formula}
\[
 \exp(x+iy) = e^x \cdot (\cos y + i \sin y).
\]
This function has the following properties:
\begin{enumerate}
  \item $\exp(z+w) = \exp z \cdot \exp w$;
  \item $\exp\bar z = \overline{\exp z}$;
  \item $\abs{\exp z} = e^{\Re z}$;
  \item $\exp x = e^x$ if $x\in \RR$.
\end{enumerate}
Since $\exp$ extends the real exponential function $e^x$, it is also natural to
use the same notation by setting:
\mynote{%
The notation $e^z$ should not lead us to think that we have defined an operation
of exponentiation between complex numbers. In fact, it is not possible to define
the power $z^w$ in a meaningful and unique way if the base $z$ is not a positive
real number. This is related to the fact that the exponential function $e^z$ is
not uniquely invertible, and thus the logarithm $\ln z$ can only be defined as a
\emph{multivalued} function.
}%
\[
  e^z = \exp z = e^x\cdot (\cos y + i\sin y), \qquad \text{if $z=x+iy$}.  
\]

Note that the complex exponential function is $2\pi i$ periodic; indeed, if
$z=x+iy$, we have
\[
\exp(z+2\pi i) 
= e^x\cdot (\cos(y+2\pi) + i\sin(y+2\pi)) 
= e^x\cdot (\cos y + i \sin y) 
= \exp(z).
\]

\subsection{Polar Representation of Complex Numbers}

As we have defined them, complex numbers $z\in \CC$ can be represented in the
form:
\[
  z = x + i y.
\]
This representation of complex numbers is called \emph{Cartesian} because it
corresponds each number $z$ to its coordinates $(x,y)$ in the complex plane
(Gauss plane).

Thanks to the definition of the complex exponential, we can also provide a
\emph{polar} representation of complex numbers. If $z=x+iy$ is any complex
number, we can define $\rho = \abs{z} = \sqrt{x^2+y^2}$ as its \emph{modulus},
which is the geometric distance between the point $z$ in the complex plane and
the origin $0\in \CC$. If $z\neq 0$, we can define the measure of the angle
formed by $z$ with the $x$-axis as the unique $\theta \in
\closeopeninterval{0}{2\pi}$ such that
\[
  \begin{cases}
    x = \rho \cos \theta,\\
    y = \rho \sin \theta.
  \end{cases}
\]
Thus, we have 
\[
  z = \rho \cdot (\cos \theta + i \sin \theta). 
\]
The pair of numbers $(\theta,\rho)$ with $\rho>0$ and
$\theta\in\closeopeninterval{0}{2\pi}$ are called the \emph{polar coordinates}
of the complex number $z$ and uniquely identify $z$. Due to the periodicity of
the functions $\sin$ and $\cos$, it is clear that the angle $\theta$ can be
replaced with $\theta +2k\pi$ for any $k\in \ZZ$ without changing the point $z$.
If $z=0$, then $\rho=\abs{z}=0$ and the coordinate $\theta$ is irrelevant.

The \emph{exponential} representation of a complex number is essentially
identical to the polar representation but uses the complex exponential instead
of trigonometric functions. Since $e^{i\theta} = \cos \theta + i\sin \theta$, if
$z=\rho\cdot (\cos \theta + i \sin \theta)$, we can write:
\[
  z = \rho \cdot e^{i\theta}.  
\]

If $z=\rho e^{i\theta}$, the number $\theta$ is usually called the
\emph{argument} of the complex number $z$ and is sometimes denoted as follows:
\[
  \theta = \arg z.  
\]
The definition of argument is inherently ambiguous since $\theta$ is not
uniquely determined (instead of $\theta$, we can choose $\theta+2k\pi)$ for any
$k\in \ZZ$). To have a unique definition, one can impose the condition
$\theta\in\closeopeninterval{0}{2\pi}$.
\mynote{But the condition 
$\theta\in\closeinterval{-\pi}{\pi}$ 
would work equally well.}
If $z=0$, we can arbitrarily define $\arg z = 0$. In formulas, we have:
\[
  \arg z =
  \begin{cases}
%   \arctg \frac y x & \text{if $x>0$,} \\
   \frac \pi 2 - \arctg \frac x y & \text{if $y>0$,} \\
   \frac 3 2 \pi- \arctg \frac x y & \text{if $y<0$,} \\
   \pi & \text{if $y=0$ and $x<0$,} \\
   0 & \text{if $y=0$ and $x\ge 0$.}
   \end{cases}
\]

\subsection{Complex $n$-th Roots}

Let $c\in \CC$ be a complex number with $c\neq 0$. We pose the problem of
determining the complex solutions of the equation
\[
  z^n = c.
\]
Such solutions will be called \emph{$n$-th roots}%
\mymargin{$n$-th roots}%
\index{root!$n$-th} of $c$.

We write $c$ and $z$ in exponential form:
\[
  c = r e^{i\alpha}, \qquad
  z = \rho e^{i\theta}.
\]
Then we have
\[
  z^n = \rho^n (e^{i\theta})^n = \rho^n e^{i n \theta}.
\]
In order for $z^n = c$, we must have equality of the moduli, i.e., $\rho^n = r$,
and equality up to integer multiples of $2\pi$ of the arguments: $n \theta =
\alpha + 2 k \pi$ with $k\in \ZZ$. Thus, we find
\[
  \theta = \frac{\alpha}{n} + k\frac{2\pi}{n}
\qquad k \in \ZZ.
\]
Now observe that for $k=0,\dots, n-1$, the second term $k 2\pi /n$ assumes $n$
distinct values within $[0,2\pi)$. For other values of $k$, we obtain angles
that differ from these by a multiple of $2\pi$, and thus no additional solutions
are found.

Therefore, the equation $z^n = c$ for $c\neq 0$ has $n$ distinct solutions given
by
\[
z_k = \sqrt[n]{r} \cdot e^{i\alpha/n + 2k\pi i /n},
\qquad k=0,1, \dots, n-1
\]
where $\alpha = \arg(c)$ and $r = \abs{c}$. From a geometric perspective, it is
observed that $z_0$ is the complex number with modulus the $n$-th root of the
given number $c$ and argument equal to one $n$-th of the argument of $c$. All
other solutions lie on the circle centered at $0$ and passing through $z_0$, and
together with $z_0$, they form the vertices of a regular $n$-gon.

In particular, in the case $c=1$, it is observed that the $n$-th roots of unity
are geometrically represented as the vertices of the regular $n$-gon inscribed
in the unit circle with one vertex at $z_0=1$.

\begin{exercise}
Find the solutions $z \in \CC$ of the following equations. Write the solutions
in polar and Cartesian form.
\begin{gather*}
   z^4 = -4 \\
   z^6 = i\\
   z^3 = -8i \\
   z^4 = z\\
   z^2 + 1 = i\sqrt{3} \\
   (z-i)^4 = 1\\
   1 + z + z^2 + z^3 = 0\\
   z^{14} - z^6 - z^8 + 1 = 0
\end{gather*}
\end{exercise}